\documentclass{article}

% if you need to pass options to natbib, use, e.g.:
%     \PassOptionsToPackage{numbers, compress}{natbib}
% before loading neurips_2025

% The authors should use one of these tracks.
% Before accepting by the NeurIPS conference, select one of the options below.
% 0. "default" for submission
 \usepackage[main, final]{neurips_2025}
% the "default" option is equal to the "main" option, which is used for the Main Track with double-blind reviewing.
% 1. "main" option is used for the Main Track
%  \usepackage[main]{neurips_2025}
% 2. "position" option is used for the Position Paper Track
%  \usepackage[position]{neurips_2025}
% 3. "dandb" option is used for the Datasets \& Benchmarks Track
 % \usepackage[dandb]{neurips_2025}
% 4. "creativeai" option is used for the Creative AI Track
%  \usepackage[creativeai]{neurips_2025}
% 5. "sglblindworkshop" option is used for the Workshop with single-blind reviewing
 % \usepackage[sglblindworkshop]{neurips_2025}
% 6. "dblblindworkshop" option is used for the Workshop with double-blind reviewing
%  \usepackage[dblblindworkshop]{neurips_2025}

% After being accepted, the authors should add "final" behind the track to compile a camera-ready version.
% 1. Main Track
 % \usepackage[main, final]{neurips_2025}
% 2. Position Paper Track
%  \usepackage[position, final]{neurips_2025}
% 3. Datasets \& Benchmarks Track
 % \usepackage[dandb, final]{neurips_2025}
% 4. Creative AI Track
%  \usepackage[creativeai, final]{neurips_2025}
% 5. Workshop with single-blind reviewing
%  \usepackage[sglblindworkshop, final]{neurips_2025}
% 6. Workshop with double-blind reviewing
%  \usepackage[dblblindworkshop, final]{neurips_2025}
% Note. For the workshop paper template, both \title{} and \workshoptitle{} are required, with the former indicating the paper title shown in the title and the latter indicating the workshop title displayed in the footnote.
% For workshops (5., 6.), the authors should add the name of the workshop, "\workshoptitle" command is used to set the workshop title.
% \workshoptitle{WORKSHOP TITLE}

% "preprint" option is used for arXiv or other preprint submissions
 % \usepackage[preprint]{neurips_2025}

% to avoid loading the natbib package, add option nonatbib:
%    \usepackage[nonatbib]{neurips_2025}

\usepackage[utf8]{inputenc} % allow utf-8 input
\usepackage[T1]{fontenc}    % use 8-bit T1 fonts
\usepackage{hyperref}       % hyperlinks
\usepackage{url}            % simple URL typesetting
\usepackage{booktabs}       % professional-quality tables
\usepackage{amsfonts}       % blackboard math symbols
\usepackage{nicefrac}       % compact symbols for 1/2, etc.
\usepackage{microtype}      % microtypography
\usepackage{xcolor}         % colors

% Note. For the workshop paper template, both \title{} and \workshoptitle{} are required, with the former indicating the paper title shown in the title and the latter indicating the workshop title displayed in the footnote. 
\title{Midterm Project Proposal: Semantically-Aware Diffusion Transformers for High-Speed, Low-Compute Image Generation}


% The \author macro works with any number of authors. There are two commands
% used to separate the names and addresses of multiple authors: \And and \AND.
%
% Using \And between authors leaves it to LaTeX to determine where to break the
% lines. Using \AND forces a line break at that point. So, if LaTeX puts 3 of 4
% authors names on the first line, and the last on the second line, try using
% \AND instead of \And before the third author name.


\author{%
  Brody L. Jordan\\
  Texas A\&M University\\
  431005357\\
  \texttt{blj@tamu.edu} \\
  % examples of more authors
  % \And
  % Coauthor \\
  % Affiliation \\
  % Address \\
  % \texttt{email} \\
  % \AND
  % Coauthor \\
  % Affiliation \\
  % Address \\
  % \texttt{email} \\
  % \And
  % Coauthor \\
  % Affiliation \\
  % Address \\
  % \texttt{email} \\
  % \And
  % Coauthor \\
  % Affiliation \\
  % Address \\
  % \texttt{email} \\
}


\begin{document}


\maketitle


\begin{abstract}
    Diffusion Transformers (DiT) are the critical components of leading high-fidelity text-to-image synthesis models, yet they are quite costly due to uniform token processing.
\end{abstract}

\section{Introduction and Motivation}

Text-to-image synthesis has undergone a revolutionary transformation with the advent of
Diffusion Transformers (DiTs). Models such as Stable Diffusion 3 \cite{esser2024scaling}
and FLUX \cite{labs2024flux} demonstrate unprecedented visual quality by replacing
traditional U-Net backbones with scalable transformer architectures. However, this
quality comes at a steep computational cost: DiTs process \textit{every} latent token
with equal intensity, regardless of semantic importance.

The quadratic complexity $O(N^2)$ of self-attention in DiTs becomes the primary
bottleneck when scaling to high resolutions. For a $1024 \times 1024$ image with
$16\times$ VAE compression, the sequence length $N = 4096$, resulting in $\sim$16M
attention operations per layer per timestep. With typical models requiring 50
denoising steps and 20+ transformer layers, total operations exceed $16 \times 10^9$
per image generation—prohibitive for real-time applications.

This survey focuses on \textbf{Spatially-Adaptive Computation} in DiTs: methods that
reduce this burden by selectively processing tokens based on their importance.
We examine 8 high-impact papers from CVPR, NeurIPS, and ICLR (2024--2025) that
attack this problem from three distinct angles: (1) architectural efficiency,
(2) token reduction strategies, and (3) hybrid representations. Our analysis
reveals that while recent works achieve impressive speedups, they fail to
dynamically align computational allocation with the prompt's semantic requirements—a
critical gap we address in our proposed research direction.

% Define the specific sub-field or problem you are investigating.
% Explain why this problem is important in the context of Multimodal Vision or 3D/Generative AI.
% Briefly state the scope of your survey.

% 4. Synthesis of Weaknesses Across the State-of-the-Art
% Despite the rapid scaling of Diffusion Transformers (DiTs) and the adoption of Flow Matching, current approaches suffer from several systematic failures that hinder their real-world utility and theoretical robustness.
% Systematic Failures and Fragile Assumptions

% Fixed Latent Resolution vs. Adaptive Compositionality: Most models (e.g., FLUX.1, SD3) assume a fixed-grid latent representation (e.g., $128 \times 128$ latents for a $1024 \times 1024$ image). In real-world scenarios, visual complexity is non-uniformly distributed. Current models waste equal compute on a "clear blue sky" as they do on a "complex human face," leading to a scalability bottleneck where high-frequency details are often blurry or structurally inconsistent despite massive parameter counts.
% Sequential Inference vs. Temporal/Semantic Parallelism: While Flow Matching simplifies the trajectory, the inference process remains fundamentally sequential across timesteps. Distillation methods like LCM or ADD try to collapse these steps, but they often "average" the manifold, losing the stochastic diversity and fine-grained control inherent in the full ODE path. The assumption that visual information must be refined through a fixed sequence of global denoising steps is a major performance inhibitor.
% The "Vicious Cycle" of Guidance (CFG): Almost all scalable models rely on Classifier-Free Guidance (CFG) to achieve prompt alignment. However, CFG requires two forward passes (conditional and unconditional) at every step, doubling the compute cost. Furthermore, high CFG scales often lead to over-saturation and artifacting, suggesting that the underlying model's "alignment" with text is shallower than the scaling laws imply.
% Language-Vision Semantic Gap: Models like SD3 and FLUX have scaled the text encoder (e.g., T5-XXL), but the interaction between text and image tokens in the DiT backbone remains relatively simplistic (concatenation or cross-attention). This assumes that high-level linguistic concepts can be mapped linearly to visual latents, which fails for complex spatial relationships, negations, or counting tasks.
\section{Categorization and Taxonomy}

We organize the surveyed literature into three complementary paradigms:

\subsection{Category A: Architectural Efficiency Improvements}
These methods modify the DiT backbone to reduce computational complexity
\textit{without} reducing token count, through attention approximations or
structural changes.

\textbf{Papers:}
\begin{itemize}
    \item \textbf{SANA 1.5} \cite{xie2025sana}: Linear Diffusion Transformers
          with depth-wise convolution attention, reducing complexity from $O(N^2)$ to
          $O(N)$ while maintaining global receptive fields.
    \item \textbf{EDiT} \cite{becker2025edit}: Efficient Diffusion Transformers
          with Linear Compressed Attention using low-rank factorization of attention
          matrices.
\end{itemize}

\subsection{Category B: Token Reduction and Merging}
These approaches explicitly reduce the sequence length $N$ by merging or
pruning redundant tokens during inference.

\textbf{Papers:}
\begin{itemize}
    \item \textbf{Token Merging (ToMe)} \cite{bolya2023token}: Bipartite
          matching-based merging for Vision Transformers, adapted for diffusion models.
    \item \textbf{ToMA} \cite{lu2025toma}: Token Merge with Attention,
          introducing saliency-aware merging strategies specifically for DiTs.
    \item \textbf{Dynamic Chunking DiT} \cite{haridas2026dynamic}: Recent work
          introducing patch-level scheduling based on local information density.
\end{itemize}

\subsection{Category C: Hybrid Representations and Alternative Paradigms}
These works challenge the pure-diffusion paradigm by combining discrete
autoregressive modeling with continuous refinement.

\textbf{Papers:}
\begin{itemize}
    \item \textbf{HART} \cite{tang2024hart}: Hybrid Autoregressive Transformer
          combining discrete tokens for structure with continuous residuals for
          detail, achieving single-pass generation.
    \item \textbf{DDiT} \cite{kim2026ddit}: Dynamic Patch Scheduling using
          learnable region importance scores.
    \item \textbf{Shiva-DiT} \cite{zhang2026shiva}: Residual-Based
          Differentiable Top-$k$ Selection for token pruning.
\end{itemize}

\begin{table}[h]
    \centering
    \caption{Taxonomy of Efficient DiT Methods}
    \label{tab:taxonomy}
    \begin{tabular}{@{}lccc@{}}
        \toprule
        \textbf{Method}  & \textbf{Category}  & \textbf{Speedup} & \textbf{Granularity}       \\
        \midrule
        SANA 1.5         & A: Architectural   & 2--3$\times$     & Global                     \\
        EDiT             & A: Architectural   & 1.5--2$\times$   & Global                     \\
        ToMe             & B: Token Reduction & 2$\times$        & Spatial (static)           \\
        ToMA             & B: Token Reduction & 2.5$\times$      & Spatial (attention-guided) \\
        Dynamic Chunking & B: Token Reduction & 3$\times$        & Patch-level                \\
        Shiva-DiT        & B: Token Reduction & 2--4$\times$     & Token-level (residual)     \\
        HART             & C: Hybrid          & 10$\times$       & Representation-level       \\
        \bottomrule
    \end{tabular}
\end{table}

% Do not just list the 6–10 papers sequentially. Group them logically.
% How do these papers approach the problem differently? (e.g., split by representation type, optimization strategy, architecture, etc.).
\section{Deep Dive and Methodological Comparison}

\subsection{Architectural Efficiency: The SANA Approach}

SANA 1.5 \cite{xie2025sana} replaces standard self-attention with
\textbf{Linear Diffusion Transformer} blocks using depth-wise convolutions.
The attention mechanism becomes:

% $\text{Attention}(Q, K, V) = \text{Softmax}\left(\frac{QK^T}{\sqrt{d}}\right)V
%     \approx \text{DWConv}(Q) \cdot V$

While this achieves linear complexity $O(N)$, it \textit{uniformly} reduces
computation across all tokens. Regions requiring fine detail (faces, text)
receive the same ``lightweight'' attention as simple backgrounds, potentially
compromising high-frequency generation.

\textbf{Trade-off:} SANA trades model expressiveness for speed. The linear
approximation cannot capture long-range dependencies as effectively as full
attention, particularly for complex compositional scenes.

\subsection{Token Reduction Strategies: From Static to Dynamic}

\paragraph{ToMe (Static Merging).}
ToMe \cite{bolya2023token} defines ``similarity'' via cosine distance in feature
space, merging the $r$ most similar token pairs at each layer:

$$
    \mathcal{M} = \{(i, j) \mid \text{argsort-top-}r(\cos(x_i, x_j))\}
$$

\textbf{Limitation:} Feature similarity does not correlate with semantic
importance. Two distinct texture patches may have high feature similarity
but require independent processing for photorealism.

\paragraph{ToMA (Attention-Guided Merging).}
ToMA \cite{lu2025toma} improves upon ToMe by using cross-attention maps
$A \in \mathbb{R}^{N \times M}$ (where $M$ is text sequence length):

$$
    s_i = \sum_{j=1}^{M} A_{i,j}
$$

Tokens with low $s_i$ (poor alignment with text) are prioritized for merging.

\textbf{Advancement:} Text-to-vision alignment improves merging decisions.
However, ToMA applies merging \textit{statically} at fixed layers, failing
to adapt to the evolving importance of tokens across denoising timesteps.

\paragraph{DDiT and Dynamic Chunking.}
Recent work \cite{kim2026ddit, haridas2026dynamic} introduces temporal
adaptivity: token importance is recomputed at each denoising step using
learned ``Decider'' modules or patch-level statistics.

\textbf{Gap:} These methods schedule compute at the \textit{patch} or
\textit{chunk} level, not individual tokens. Coarse-grained decisions
miss opportunities for fine-grained efficiency.

\paragraph{Shiva-DiT.}
Shiva-DiT \cite{zhang2026shiva} implements differentiable Top-$k$ selection
using the Gumbel-Softmax trick:

$$
    P(\text{keep}_i) = \text{Softmax}\left(\frac{\log \pi_i + g_i}{\tau}\right)
$$

where $\pi_i$ is a learned importance score and $g_i \sim \text{Gumbel}(0,1)$.

\textbf{Innovation:} End-to-end trainable sparsity. However, the importance
score $\pi_i$ is learned from reconstruction loss alone, lacking explicit
\textit{text-prompt guidance}.

\subsection{Hybrid Representations: HART}

HART \cite{tang2024hart} abandons the diffusion paradigm entirely for a
hybrid autoregressive approach:

\begin{enumerate}
    \item Discrete VQ-VAE tokens capture global structure in one AR pass.
    \item Continuous residuals add detail via small diffusion heads.
\end{enumerate}

\textbf{Strengths:} Achieves single-pass generation (massive speedup),
leveraging the efficiency of AR models (GPT-style next-token prediction).

\textbf{Fundamental Limitation:} HART (and all AR models) require
\textit{discrete tokenization}, which introduces an information bottleneck
at the VQ-VAE stage. Complex textures and continuous gradients suffer from
quantization artifacts. Furthermore, AR models autoregressively decode,
making them inherently sequential—batch generation of multiple images
cannot be parallelized as effectively as diffusion.

% This is the core of your survey. Critically compare the methodologies of the surveyed papers.
% What are the trade-offs between these approaches? (e.g., rendering speed vs. memory footprint, visual quality vs. training time).
% Include technical depth; do not shy away from discussing the underlying math, loss functions, or architectural choices.
\section{Limitations and Open Challenges}

Our survey reveals three systematic failures across the state-of-the-art:

\subsection{Challenge 1: The Uniformity Assumption}
\textbf{Current approaches assume uniform importance across the spatial
    domain.} Whether applying full attention (standard DiT), linear attention
(SANA), or token merging (ToMe), existing methods process the ``sky''
region with the same intensity as a ``human face.''

\textbf{Why it fails:} Visual scenes exhibit extreme spatial redundancy.
In most images, 70--80\% of pixels are low-entropy (smooth regions,
repetitive textures). Current architectures waste compute on these
``uninformative'' regions while under-investing in high-frequency details
critical for perceived quality.

\subsection{Challenge 2: Static vs. Dynamic Importance}
\textbf{Text-prompt alignment is underutilized in compute allocation.}
While ToMA uses cross-attention for merging decisions, it treats this
saliency as \textit{constant} across timesteps.

\textbf{Why it fails:} The importance of a token evolves during
denoising. At $t=T$ (pure noise), all tokens are equally uncertain.
At $t=T/2$, some regions (e.g., backgrounds) may stabilize while others
(e.g., foreground objects) remain noisy. Current methods lack mechanisms
to dynamically re-prioritize tokens as the diffusion trajectory unfolds.

\subsection{Challenge 3: The Differentiable Trilemma}
\textbf{End-to-end trainable sparsity requires balancing three conflicting
    goals:} (1) differentiability (for gradient-based training), (2) efficiency
(for wall-clock speedup), and (3) quality preservation (no generation
artifacts).

Shiva-DiT achieves (1) and (2) but compromises on (3)—learned importance
scores drift toward conservative selection, reducing effective sparsity.
Feature caching methods achieve (2) and (3) but are \textit{not}
differentiable, requiring hand-tuned heuristics.

\textbf{Why it matters:} Without a unified solution to this trilemma,
spatial sparsity remains a ``post-hoc'' optimization rather than a
first-class architectural primitive.

\subsection{Challenge 4: Seamless Integration}
\textbf{No existing method handles the boundary between sparse and dense
    regions.} When merging tokens or dropping computation, current approaches
risk visible artifacts at region boundaries—the ``blocky'' effect where
high-detail subjects meet low-detail backgrounds.

\textbf{Real-world impact:} These limitations prevent deployment of
high-quality DiTs on consumer hardware. Current models require 24GB+ VRAM
and 10--30 seconds per image, limiting applications in real-time creative
tools, mobile AR, and autonomous systems requiring rapid visual
imagination.

% This is the most important section for your grade. Synthesize the weaknesses across the state-of-the-art.
% What are the systematic failures of the current approaches? What assumptions do they make that do not hold in real-world scenarios?
% Your critique should demonstrate a deep enough understanding to justify why a new NeurIPS-level solution is needed.
\section{Proposed Research Direction}
\subsection{The Specific Limitation to Solve}
We target \textbf{Challenge 2 (Dynamic Importance)} and \textbf{Challenge 4
    (Seamless Integration)}. Our hypothesis is that \textit{compute allocation
    should be a function of both the text prompt (semantic saliency) and the
    current timestep (temporal stability).}

\subsection{Proposed Solution: SA-SparseDiT}

We propose \textbf{Spatially-Adaptive Sparse Diffusion Transformers
    (SA-SparseDiT)}, a DiT variant with three novel components:

\subsubsection{1. Tri-Signal Saliency Router}
A differentiable gating module that computes per-token priority $P_i$ using
three signals:

$$P_i = \underbrace{\alpha \cdot A_{\text{text}, i}}_{\text{Text Alignment}} +
    \underbrace{\beta \cdot \|\nabla_\theta \epsilon_i\|_2}_{\text{Noise Magnitude}} +
    \underbrace{\gamma \cdot \mathcal{V}(x_i)}_{\text{Spatial Variance}}$$

where:
\begin{itemize}
    \item $A_{\text{text}, i}$: Cross-attention weight to text embeddings
          (semantic importance).
    \item $\|\nabla_\epsilon\|_2$: Magnitude of predicted noise (temporal
          uncertainty).
    \item $\mathcal{V}(x_i)$: Local spatial variance (structural complexity).
\end{itemize}

\textbf{Why this is novel:} Unlike ToMA (static attention) or Shiva-DiT
(residual only), our router dynamically recomputes priorities at each
denoising step, allowing tokens to ``graduate'' from sparse to dense
processing as they become semantically relevant.

\subsubsection{2. Differentiable Top-$k$ Selection via Gumbel-Softmax}
Using the concrete distribution \cite{jang2016categorical}, we implement
sparse selection:

$$
    z_i = \text{sigmoid}\left(\frac{\log P_i + g_i}{\tau}\right), \quad g_i \sim \text{Gumbel}(0,1)
$$

\textbf{Top-$k$ tokens} ($z_i > \tau_{\text{sparsity}}$) proceed through
full self-attention. \textbf{Bottom-$(N-k)$ tokens} are merged into $m \ll k$
anchor tokens via a learned clustering operation.

\subsubsection{3. Seamless Re-Integration via Residual Flows}
To address Challenge 4, we introduce \textbf{seamless un-pooling} using
normalizing flows \cite{dinh2016density}. When expanding merged tokens
back to full resolution, we apply invertible transformations that
preserve local smoothness while allowing high-frequency detail injection
from dense regions.

\subsection{Expected Contributions}
1. \textbf{First prompt-driven dynamic sparsity} for DiTs with per-token,
per-timestep re-prioritization.
2. \textbf{Differentiable sparse attention kernels} enabling end-to-end
training without quality degradation.
3. **3--5$\times$ speedup** on 1024p generation with $<5\%$ FID increase,
enabling 8K synthesis on consumer GPUs.

\subsection{Risk Mitigation}
\begin{itemize}
    \item \textbf{Rebuttal: ``HART already solves efficiency.''}
          $\rightarrow$ HART requires discrete tokenization. SA-SparseDiT
          preserves the continuous flow-matching paradigm, maintaining superior
          texture quality.
    \item \textbf{Rebuttal: ``Dynamic Chunking DDiT already token-prunes.''}
          $\rightarrow$ DDiT operates at patch-level ($16\times16$ blocks).
          SA-SparseDiT is token-level, offering finer granularity.
\end{itemize}

\begin{thebibliography}{99}

    \bibitem{esser2024scaling}
    Patrick Esser et al.
    \newblock Scaling Rectified Flow Transformers for High-Resolution Image Synthesis.
    \newblock \emph{arXiv preprint arXiv:2403.03206}, 2024.

    \bibitem{labs2024flux}
    Black Forest Labs.
    \newblock FLUX.1: A Series of State-of-the-Art Diffusion Models.
    \newblock \emph{Technical Report}, 2024.

    \bibitem{xie2025sana}
    Enze Xie et al.
    \newblock SANA 1.5: Efficient Scaling of Training-Time and Inference-Time Compute in Linear Diffusion Transformer.
    \newblock \emph{arXiv preprint arXiv:2501.18427}, 2025.

    \bibitem{becker2025edit}
    Philipp Becker et al.
    \newblock EDiT: Efficient Diffusion Transformers with Linear Compressed Attention.
    \newblock \emph{arXiv preprint arXiv:2503.16726}, 2025.

    \bibitem{bolya2023token}
    Daniel Bolya and Judy Hoffman.
    \newblock Token Merging for Fast Stable Diffusion.
    \newblock \emph{arXiv preprint arXiv:2303.17604}, 2023.

    \bibitem{lu2025toma}
    Wenbo Lu et al.
    \newblock ToMA: Token Merge with Attention for Diffusion Models.
    \newblock \emph{arXiv preprint arXiv:2509.10918}, 2025.

    \bibitem{haridas2026dynamic}
    Akash Haridas et al.
    \newblock Dynamic Chunking Diffusion Transformer.
    \newblock \emph{arXiv preprint arXiv:2603.06351}, 2026.

    \bibitem{tang2024hart}
    Haotian Tang et al.
    \newblock HART: Efficient Visual Generation with Hybrid Autoregressive Transformer.
    \newblock \emph{arXiv preprint arXiv:2410.10812}, 2024.

    \bibitem{kim2026ddit}
    Dahye Kim et al.
    \newblock DDiT: Dynamic Patch Scheduling for Efficient Diffusion Transformers.
    \newblock \emph{arXiv preprint arXiv:2602.16968}, 2026.

    \bibitem{zhang2026shiva}
    Jiaji Zhang et al.
    \newblock Shiva-DiT: Residual-Based Differentiable Top-k Selection for Efficient Diffusion Transformers.
    \newblock \emph{arXiv preprint arXiv:2602.05605}, 2026.

    \bibitem{jang2016categorical}
    Eric Jang, Shixiang Gu, and Ben Poole.
    \newblock Categorical Reparameterization with Gumbel-Softmax.
    \newblock \emph{ICLR}, 2017.

    \bibitem{dinh2016density}
    Laurent Dinh, Jascha Sohl-Dickstein, and Samy Bengio.
    \newblock Density estimation using Real NVP.
    \newblock \emph{ICLR}, 2017.

\end{thebibliography}

% Proposed Research Direction (Towards a NeurIPS Submission)
% Specific Limitation to Solve: The Computational Inefficiency of Uniform Processing in high-resolution synthesis. Currently, $O(N^2)$ attention in DiTs (or even $O(N)$ linear approximations) processes all patches with equal priority, regardless of their semantic density or relevance to the prompt.
% Proposed Direction: Spatially-Adaptive Sparse Diffusion Transformers (SA-SparseDiT).
% Technical Intuition and Hypothesis:
% In natural images, semantic information is sparse. A NeurIPS-level solution should move away from dense, grid-based processing. My initial hypothesis is that we can formulate the denoising process as a Sparse-to-Dense Token Evolution.

% Initial Intuition: Instead of processing all latent tokens at every timestep, the model should use a "Saliency-Aware Router" (potentially inspired by MoE or top-$k$ selection) to identify tokens requiring high-fidelity refinement based on the text prompt.
% Technical Approach:

% Dynamic Token Pruning/Merging: During the early timesteps (coarse structure), the model operates on a heavily downsampled or merged set of tokens.
% Prompt-Driven Saliency: As the trajectory progresses, the model uses the cross-attention maps between the text and image to "un-pool" or "reactivate" specific spatial regions (e.g., a person's eyes or a complex texture) for dense refinement, while keeping background regions sparse.
% Hypothesis: By allocating compute dynamically based on spatial complexity and prompt relevance, we can achieve $2\times$ to $4\times$ inference speedups and better high-resolution detail without the quality loss typical of global distillation. This shifts the scaling law from "more parameters for better pixels" to "better compute allocation for relevant pixels." \\

% The proposed research direction, Spatially-Adaptive Sparse Diffusion Transformers (SA-SparseDiT), shifts the fundamental paradigm of image generation from dense, uniform computation to sparse, importance-based allocation.
% In current state-of-the-art models like FLUX.1 or Stable Diffusion 3, every single pixel (or latent token) is treated with equal importance. Whether the model is denoising a high-contrast human eye or a flat, out-of-focus background, it executes the same number of floating-point operations (FLOPs). This is computationally wasteful and mathematically inefficient.
% Here is a detailed breakdown of the technical intuition, the architectural components, and the reasoning behind this direction. \\

% \subsection{The Core Hypothesis}
% 1. The Core Hypothesis: Semantic Sparsity
% The fundamental intuition is that visual information density is non-uniform.

% Early Stages (Structural): At high noise levels (early timesteps), the model only needs to establish global layout and color masses. High-resolution, token-by-token processing is redundant here.
% Late Stages (Refinement): At low noise levels, the model needs high resolution, but only in "semantically rich" areas (e.g., faces, text, intricate patterns). The "background" often reaches its final state much earlier than the "subject."

% Hypothesis: If we can dynamically identify which tokens "matter" at each step of the ODE trajectory, we can prune or merge the "unimportant" tokens to drastically reduce the $O(N^2)$ or even $O(N)$ cost of the Transformer blocks.

% \subsection{Architectural Components}
% 2. Architectural Components
% A. The Prompt-Driven Saliency Router
% Instead of a static grid, the model uses the Cross-Attention maps from the first few layers to generate a "Saliency Mask."

% If the prompt is "A cat sitting on a wooden floor," the router identifies that tokens corresponding to the "cat" and the "wood grain" require high-frequency updates.
% The "floor" tokens in the distant background are marked for "Sparse Processing."

% B. Dynamic Token Merging \& Splitting
% Inspired by methods like Token Merging (ToMe) but made adaptive and reversible:

% Merging: In "low-importance" regions, similar tokens are clustered into a single "Representative Token." This reduces the sequence length $L$ that the Self-Attention mechanism must process.
% Splitting (Un-pooling): As the denoising progresses, if the model detects a sudden increase in local variance or a new prompt-match in a sparse region, it "splits" the representative token back into high-resolution patches for fine-detail synthesis.

% C. The Multi-Resolution Trajectory
% The model doesn't just stay in one latent space. It follows a V-shaped compute trajectory:

% Phase 1 (Global): All tokens are merged/downsampled. The model defines the "where" of the objects.
% Phase 2 (Sparse-Dense): Tokens corresponding to the prompt-subject are un-merged to full resolution; background tokens remain merged.
% Phase 3 (Cleanup): A final high-speed pass over the entire grid to ensure global consistency.

% \subsection{NeurIPS}
% Why this is a NeurIPS-level Solution
% Current "efficiency" papers usually focus on Distillation (reducing the number of steps) or Linear Attention (changing the math of the layer). SA-SparseDiT is more ambitious because it changes the data flow itself.






























% FeatureStandard DiT (e.g., SD3/FLUX)Proposed SA-SparseDiTToken ProcessingDense \& Uniform (All tokens, all steps)Sparse \& Adaptive (Importance-weighted)Complexity$O(T \times L^2)$ or $O(T \times L)$ $O(T \times Active Tokens)$ Information Theory Ignores spatial redundancy Exploits spatial and temporal redundancyPrompt AlignmentPassive (Text as bias)Active (Text drives compute allocation)
% Technical Challenges (The "Research" part)

% Differentiable Routing: How do you train a router to pick tokens without breaking the gradient flow during training? (Possible solution: Gumbel-Softmax or Reinforcement Learning).
% Artifact Prevention: Merging tokens can cause "seams" in the image. The research would involve designing a "Seamless Re-pooling" layer that uses neighborhood context to smooth boundaries between sparse and dense regions.
% Hardware Compatibility: Standard GPUs love dense matrices. Making "sparse" operations actually faster in PyTorch requires custom kernels (e.g., using Triton) to ensure the theoretical speedup translates to wall-clock time.

% \subsection{Expected Impact}
% By solving the uniformity bottleneck, this research would allow us to generate 8K or 16K images on consumer hardware. While current models "choke" as resolution increases because they try to attend to every pixel, SA-SparseDiT would only increase compute for the complex parts of the 16K image, making extreme-high-resolution synthesis a reality.
% \subsubsection{Summary}
% "We are currently wasting 70-80\% of our generative compute on 'background noise' and redundant pixels.
% By introducing SA-SparseDiT, we align the computational structure of the model with the semantic structure of the human visual
% system. This isn't just an 'optimization trick'--it's a fundamental architectural shift that enables the next generation of
% high-resolution, long-form, and on-device visual intelligence."

% Based directly on the limitations identified in Section 4, propose a concrete direction for your own research (which will ideally evolve into your final project for this course).
% What specific limitation are you going to solve, and what is your initial hypothesis or technical intuition for solving it?
\pagebreak



% \section{Submission of papers to NeurIPS 2025}


% Please read the instructions below carefully and follow them faithfully.


% \subsection{Style}


% Papers to be submitted to NeurIPS 2025 must be prepared according to the
% instructions presented here. Papers may only be up to {\bf nine} pages long,
% including figures.
% % Additional pages \emph{containing only acknowledgments and references} are allowed.
% Additional pages \emph{containing references, checklist, and the optional technical appendices} do not count as content pages.
% Papers that exceed the page limit will not be
% reviewed, or in any other way considered for presentation at the conference.


% The margins in 2025 are the same as those in previous years.


% Authors are required to use the NeurIPS \LaTeX{} style files obtainable at the
% NeurIPS website as indicated below. Please make sure you use the current files
% and not previous versions. Tweaking the style files may be grounds for
% rejection.


% \subsection{Retrieval of style files}


% The style files for NeurIPS and other conference information are available on
% the website at
% \begin{center}
%     \url{https://neurips.cc}
% \end{center}
% The file \verb+neurips_2025.pdf+ contains these instructions and illustrates the
% various formatting requirements your NeurIPS paper must satisfy.


% The only supported style file for NeurIPS 2025 is \verb+neurips_2025.sty+,
% rewritten for \LaTeXe{}.  \textbf{Previous style files for \LaTeX{} 2.09,
%     Microsoft Word, and RTF are no longer supported!}


% The \LaTeX{} style file contains three optional arguments: \verb+final+, which
% creates a camera-ready copy, \verb+preprint+, which creates a preprint for
% submission to, e.g., arXiv, and \verb+nonatbib+, which will not load the
% \verb+natbib+ package for you in case of package clash.


% \paragraph{Preprint option}
% If you wish to post a preprint of your work online, e.g., on arXiv, using the
% NeurIPS style, please use the \verb+preprint+ option. This will create a
% nonanonymized version of your work with the text ``Preprint. Work in progress.''
% in the footer. This version may be distributed as you see fit, as long as you do not say which conference it was submitted to. Please \textbf{do
%     not} use the \verb+final+ option, which should \textbf{only} be used for
% papers accepted to NeurIPS.


% At submission time, please omit the \verb+final+ and \verb+preprint+
% options. This will anonymize your submission and add line numbers to aid
% review. Please do \emph{not} refer to these line numbers in your paper as they
% will be removed during generation of camera-ready copies.


% The file \verb+neurips_2025.tex+ may be used as a ``shell'' for writing your
% paper. All you have to do is replace the author, title, abstract, and text of
% the paper with your own.


% The formatting instructions contained in these style files are summarized in
% Sections \ref{gen_inst}, \ref{headings}, and \ref{others} below.


% \section{General formatting instructions}
% \label{gen_inst}


% The text must be confined within a rectangle 5.5~inches (33~picas) wide and
% 9~inches (54~picas) long. The left margin is 1.5~inch (9~picas).  Use 10~point
% type with a vertical spacing (leading) of 11~points.  Times New Roman is the
% preferred typeface throughout, and will be selected for you by default.
% Paragraphs are separated by \nicefrac{1}{2}~line space (5.5 points), with no
% indentation.


% The paper title should be 17~point, initial caps/lower case, bold, centered
% between two horizontal rules. The top rule should be 4~points thick and the
% bottom rule should be 1~point thick. Allow \nicefrac{1}{4}~inch space above and
% below the title to rules. All pages should start at 1~inch (6~picas) from the
% top of the page.


% For the final version, authors' names are set in boldface, and each name is
% centered above the corresponding address. The lead author's name is to be listed
% first (left-most), and the co-authors' names (if different address) are set to
% follow. If there is only one co-author, list both author and co-author side by
% side.


% Please pay special attention to the instructions in Section \ref{others}
% regarding figures, tables, acknowledgments, and references.

% \section{Headings: first level}
% \label{headings}


% All headings should be lower case (except for first word and proper nouns),
% flush left, and bold.


% First-level headings should be in 12-point type.


% \subsection{Headings: second level}


% Second-level headings should be in 10-point type.


% \subsubsection{Headings: third level}


% Third-level headings should be in 10-point type.


% \paragraph{Paragraphs}


% There is also a \verb+\paragraph+ command available, which sets the heading in
% bold, flush left, and inline with the text, with the heading followed by 1\,em
% of space.


% \section{Citations, figures, tables, references}
% \label{others}


% These instructions apply to everyone.


% \subsection{Citations within the text}


% The \verb+natbib+ package will be loaded for you by default.  Citations may be
% author/year or numeric, as long as you maintain internal consistency.  As to the
% format of the references themselves, any style is acceptable as long as it is
% used consistently.


% The documentation for \verb+natbib+ may be found at
% \begin{center}
%     \url{http://mirrors.ctan.org/macros/latex/contrib/natbib/natnotes.pdf}
% \end{center}
% Of note is the command \verb+\citet+, which produces citations appropriate for
% use in inline text.  For example,
% \begin{verbatim}
%    \citet{hasselmo} investigated\dots
% \end{verbatim}
% produces
% \begin{quote}
%     Hasselmo, et al.\ (1995) investigated\dots
% \end{quote}


% If you wish to load the \verb+natbib+ package with options, you may add the
% following before loading the \verb+neurips_2025+ package:
% \begin{verbatim}
%    \PassOptionsToPackage{options}{natbib}
% \end{verbatim}


% If \verb+natbib+ clashes with another package you load, you can add the optional
% argument \verb+nonatbib+ when loading the style file:
% \begin{verbatim}
%    \usepackage[nonatbib]{neurips_2025}
% \end{verbatim}


% As submission is double blind, refer to your own published work in the third
% person. That is, use ``In the previous work of Jones et al.\ [4],'' not ``In our
% previous work [4].'' If you cite your other papers that are not widely available
% (e.g., a journal paper under review), use anonymous author names in the
% citation, e.g., an author of the form ``A.\ Anonymous'' and include a copy of the anonymized paper in the supplementary material.


% \subsection{Footnotes}


% Footnotes should be used sparingly.  If you do require a footnote, indicate
% footnotes with a number\footnote{Sample of the first footnote.} in the
% text. Place the footnotes at the bottom of the page on which they appear.
% Precede the footnote with a horizontal rule of 2~inches (12~picas).


% Note that footnotes are properly typeset \emph{after} punctuation
% marks.\footnote{As in this example.}


% \subsection{Figures}


% \begin{figure}
%     \centering
%     \fbox{\rule[-.5cm]{0cm}{4cm} \rule[-.5cm]{4cm}{0cm}}
%     \caption{Sample figure caption.}
% \end{figure}


% All artwork must be neat, clean, and legible. Lines should be dark enough for
% purposes of reproduction. The figure number and caption always appear after the
% figure. Place one line space before the figure caption and one line space after
% the figure. The figure caption should be lower case (except for first word and
% proper nouns); figures are numbered consecutively.


% You may use color figures.  However, it is best for the figure captions and the
% paper body to be legible if the paper is printed in either black/white or in
% color.


% \subsection{Tables}


% All tables must be centered, neat, clean and legible.  The table number and
% title always appear before the table.  See Table~\ref{sample-table}.


% Place one line space before the table title, one line space after the
% table title, and one line space after the table. The table title must
% be lower case (except for first word and proper nouns); tables are
% numbered consecutively.


% Note that publication-quality tables \emph{do not contain vertical rules.} We
% strongly suggest the use of the \verb+booktabs+ package, which allows for
% typesetting high-quality, professional tables:
% \begin{center}
%     \url{https://www.ctan.org/pkg/booktabs}
% \end{center}
% This package was used to typeset Table~\ref{sample-table}.


% \begin{table}
%     \caption{Sample table title}
%     \label{sample-table}
%     \centering
%     \begin{tabular}{lll}
%         \toprule
%         \multicolumn{2}{c}{Part}                     \\
%         \cmidrule(r){1-2}
%         Name     \& Description     \& Size ($\mu$m) \\
%         \midrule
%         Dendrite \& Input terminal  \& $\sim$100     \\
%         Axon     \& Output terminal \& $\sim$10      \\
%         Soma     \& Cell body       \& up to $10^6$  \\
%         \bottomrule
%     \end{tabular}
% \end{table}

% \subsection{Math}
% Note that display math in bare TeX commands will not create correct line numbers for submission. Please use LaTeX (or AMSTeX) commands for unnumbered display math. (You really shouldn't be using \$\$ anyway; see \url{https://tex.stackexchange.com/questions/503/why-is-preferable-to} and \url{https://tex.stackexchange.com/questions/40492/what-are-the-differences-between-align-equation-and-displaymath} for more information.)

% \subsection{Final instructions}

% Do not change any aspects of the formatting parameters in the style files.  In
% particular, do not modify the width or length of the rectangle the text should
% fit into, and do not change font sizes (except perhaps in the
% \textbf{References} section; see below). Please note that pages should be
% numbered.


% \section{Preparing PDF files}


% Please prepare submission files with paper size ``US Letter,'' and not, for
% example, ``A4.''


% Fonts were the main cause of problems in the past years. Your PDF file must only
% contain Type 1 or Embedded TrueType fonts. Here are a few instructions to
% achieve this.


% \begin{itemize}


%     \item You should directly generate PDF files using \verb+pdflatex+.


%     \item You can check which fonts a PDF files uses.  In Acrobat Reader, select the
%           menu Files$>$Document Properties$>$Fonts and select Show All Fonts. You can
%           also use the program \verb+pdffonts+ which comes with \verb+xpdf+ and is
%           available out-of-the-box on most Linux machines.


%     \item \verb+xfig+ "patterned" shapes are implemented with bitmap fonts.  Use
%           "solid" shapes instead.


%     \item The \verb+\bbold+ package almost always uses bitmap fonts.  You should use
%           the equivalent AMS Fonts:
%           \begin{verbatim}
%    \usepackage{amsfonts}
% \end{verbatim}
%           followed by, e.g., \verb+\mathbb{R}+, \verb+\mathbb{N}+, or \verb+\mathbb{C}+
%           for $\mathbb{R}$, $\mathbb{N}$ or $\mathbb{C}$.  You can also use the following
%           workaround for reals, natural and complex:
%           \begin{verbatim}
%    \newcommand{\RR}{I\!\!R} %real numbers
%    \newcommand{\Nat}{I\!\!N} %natural numbers
%    \newcommand{\CC}{I\!\!\!\!C} %complex numbers
% \end{verbatim}
%           Note that \verb+amsfonts+ is automatically loaded by the \verb+amssymb+ package.


% \end{itemize}


% If your file contains type 3 fonts or non embedded TrueType fonts, we will ask
% you to fix it.


% \subsection{Margins in \LaTeX{}}


% Most of the margin problems come from figures positioned by hand using
% \verb+\special+ or other commands. We suggest using the command
% \verb+\includegraphics+ from the \verb+graphicx+ package. Always specify the
% figure width as a multiple of the line width as in the example below:
% \begin{verbatim}
%    \usepackage[pdftex]{graphicx} ...
%    \includegraphics[width=0.8\linewidth]{myfile.pdf}
% \end{verbatim}
% See Section 4.4 in the graphics bundle documentation
% (\url{http://mirrors.ctan.org/macros/latex/required/graphics/grfguide.pdf})


% A number of width problems arise when \LaTeX{} cannot properly hyphenate a
% line. Please give LaTeX hyphenation hints using the \verb+\-+ command when
% necessary.

% \begin{ack}
%     Use unnumbered first level headings for the acknowledgments. All acknowledgments
%     go at the end of the paper before the list of references. Moreover, you are required to declare
%     funding (financial activities supporting the submitted work) and competing interests (related financial activities outside the submitted work).
%     More information about this disclosure can be found at: \url{https://neurips.cc/Conferences/2025/PaperInformation/FundingDisclosure}.


%     Do {\bf not} include this section in the anonymized submission, only in the final paper. You can use the \texttt{ack} environment provided in the style file to automatically hide this section in the anonymized submission.
% \end{ack}

% \section*{References}


% References follow the acknowledgments in the camera-ready paper. Use unnumbered first-level heading for
% the references. Any choice of citation style is acceptable as long as you are
% consistent. It is permissible to reduce the font size to \verb+small+ (9 point)
% when listing the references.
% Note that the Reference section does not count towards the page limit.
% \medskip


% {
% \small


% [1] Alexander, J.A.\ \& Mozer, M.C.\ (1995) Template-based algorithms for
% connectionist rule extraction. In G.\ Tesauro, D.S.\ Touretzky and T.K.\ Leen
% (eds.), {\it Advances in Neural Information Processing Systems 7},
% pp.\ 609--616. Cambridge, MA: MIT Press.


%     [2] Bower, J.M.\ \& Beeman, D.\ (1995) {\it The Book of GENESIS: Exploring
%         Realistic Neural Models with the GEneral NEural SImulation System.}  New York:
% TELOS/Springer--Verlag.


% [3] Hasselmo, M.E., Schnell, E.\ \& Barkai, E.\ (1995) Dynamics of learning and
% recall at excitatory recurrent synapses and cholinergic modulation in rat
% hippocampal region CA3. {\it Journal of Neuroscience} {\bf 15}(7):5249-5262.
% }


% %%%%%%%%%%%%%%%%%%%%%%%%%%%%%%%%%%%%%%%%%%%%%%%%%%%%%%%%%%%%

% \appendix

% \section{Technical Appendices and Supplementary Material}
% Technical appendices with additional results, figures, graphs and proofs may be submitted with the paper submission before the full submission deadline (see above), or as a separate PDF in the ZIP file below before the supplementary material deadline. There is no page limit for the technical appendices.

% %%%%%%%%%%%%%%%%%%%%%%%%%%%%%%%%%%%%%%%%%%%%%%%%%%%%%%%%%%%%

% \newpage
% \section*{NeurIPS Paper Checklist}

% %%% BEGIN INSTRUCTIONS %%%
% The checklist is designed to encourage best practices for responsible machine learning research, addressing issues of reproducibility, transparency, research ethics, and societal impact. Do not remove the checklist: {\bf The papers not including the checklist will be desk rejected.} The checklist should follow the references and follow the (optional) supplemental material.  The checklist does NOT count towards the page
% limit.

% Please read the checklist guidelines carefully for information on how to answer these questions. For each question in the checklist:
% \begin{itemize}
%     \item You should answer \answerYes{}, \answerNo{}, or \answerNA{}.
%     \item \answerNA{} means either that the question is Not Applicable for that particular paper or the relevant information is Not Available.
%     \item Please provide a short (1-2 sentence) justification right after your answer (even for NA).
%           % \item {\bf The papers not including the checklist will be desk rejected.}
% \end{itemize}

% {\bf The checklist answers are an integral part of your paper submission.} They are visible to the reviewers, area chairs, senior area chairs, and ethics reviewers. You will be asked to also include it (after eventual revisions) with the final version of your paper, and its final version will be published with the paper.

% The reviewers of your paper will be asked to use the checklist as one of the factors in their evaluation. While "\answerYes{}" is generally preferable to "\answerNo{}", it is perfectly acceptable to answer "\answerNo{}" provided a proper justification is given (e.g., "error bars are not reported because it would be too computationally expensive" or "we were unable to find the license for the dataset we used"). In general, answering "\answerNo{}" or "\answerNA{}" is not grounds for rejection. While the questions are phrased in a binary way, we acknowledge that the true answer is often more nuanced, so please just use your best judgment and write a justification to elaborate. All supporting evidence can appear either in the main paper or the supplemental material, provided in appendix. If you answer \answerYes{} to a question, in the justification please point to the section(s) where related material for the question can be found.

% IMPORTANT, please:
% \begin{itemize}
%     \item {\bf Delete this instruction block, but keep the section heading ``NeurIPS Paper Checklist"},
%     \item  {\bf Keep the checklist subsection headings, questions/answers and guidelines below.}
%     \item {\bf Do not modify the questions and only use the provided macros for your answers}.
% \end{itemize}


% %%% END INSTRUCTIONS %%%


% \begin{enumerate}

%     \item {\bf Claims}
%     \item[] Question: Do the main claims made in the abstract and introduction accurately reflect the paper's contributions and scope?
%     \item[] Answer: \answerTODO{} % Replace by \answerYes{}, \answerNo{}, or \answerNA{}.
%     \item[] Justification: \justificationTODO{}
%     \item[] Guidelines:
%           \begin{itemize}
%               \item The answer NA means that the abstract and introduction do not include the claims made in the paper.
%               \item The abstract and/or introduction should clearly state the claims made, including the contributions made in the paper and important assumptions and limitations. A No or NA answer to this question will not be perceived well by the reviewers.
%               \item The claims made should match theoretical and experimental results, and reflect how much the results can be expected to generalize to other settings.
%               \item It is fine to include aspirational goals as motivation as long as it is clear that these goals are not attained by the paper.
%           \end{itemize}

%     \item {\bf Limitations}
%     \item[] Question: Does the paper discuss the limitations of the work performed by the authors?
%     \item[] Answer: \answerTODO{} % Replace by \answerYes{}, \answerNo{}, or \answerNA{}.
%     \item[] Justification: \justificationTODO{}
%     \item[] Guidelines:
%           \begin{itemize}
%               \item The answer NA means that the paper has no limitation while the answer No means that the paper has limitations, but those are not discussed in the paper.
%               \item The authors are encouraged to create a separate "Limitations" section in their paper.
%               \item The paper should point out any strong assumptions and how robust the results are to violations of these assumptions (e.g., independence assumptions, noiseless settings, model well-specification, asymptotic approximations only holding locally). The authors should reflect on how these assumptions might be violated in practice and what the implications would be.
%               \item The authors should reflect on the scope of the claims made, e.g., if the approach was only tested on a few datasets or with a few runs. In general, empirical results often depend on implicit assumptions, which should be articulated.
%               \item The authors should reflect on the factors that influence the performance of the approach. For example, a facial recognition algorithm may perform poorly when image resolution is low or images are taken in low lighting. Or a speech-to-text system might not be used reliably to provide closed captions for online lectures because it fails to handle technical jargon.
%               \item The authors should discuss the computational efficiency of the proposed algorithms and how they scale with dataset size.
%               \item If applicable, the authors should discuss possible limitations of their approach to address problems of privacy and fairness.
%               \item While the authors might fear that complete honesty about limitations might be used by reviewers as grounds for rejection, a worse outcome might be that reviewers discover limitations that aren't acknowledged in the paper. The authors should use their best judgment and recognize that individual actions in favor of transparency play an important role in developing norms that preserve the integrity of the community. Reviewers will be specifically instructed to not penalize honesty concerning limitations.
%           \end{itemize}

%     \item {\bf Theory assumptions and proofs}
%     \item[] Question: For each theoretical result, does the paper provide the full set of assumptions and a complete (and correct) proof?
%     \item[] Answer: \answerTODO{} % Replace by \answerYes{}, \answerNo{}, or \answerNA{}.
%     \item[] Justification: \justificationTODO{}
%     \item[] Guidelines:
%           \begin{itemize}
%               \item The answer NA means that the paper does not include theoretical results.
%               \item All the theorems, formulas, and proofs in the paper should be numbered and cross-referenced.
%               \item All assumptions should be clearly stated or referenced in the statement of any theorems.
%               \item The proofs can either appear in the main paper or the supplemental material, but if they appear in the supplemental material, the authors are encouraged to provide a short proof sketch to provide intuition.
%               \item Inversely, any informal proof provided in the core of the paper should be complemented by formal proofs provided in appendix or supplemental material.
%               \item Theorems and Lemmas that the proof relies upon should be properly referenced.
%           \end{itemize}

%     \item {\bf Experimental result reproducibility}
%     \item[] Question: Does the paper fully disclose all the information needed to reproduce the main experimental results of the paper to the extent that it affects the main claims and/or conclusions of the paper (regardless of whether the code and data are provided or not)?
%     \item[] Answer: \answerTODO{} % Replace by \answerYes{}, \answerNo{}, or \answerNA{}.
%     \item[] Justification: \justificationTODO{}
%     \item[] Guidelines:
%           \begin{itemize}
%               \item The answer NA means that the paper does not include experiments.
%               \item If the paper includes experiments, a No answer to this question will not be perceived well by the reviewers: Making the paper reproducible is important, regardless of whether the code and data are provided or not.
%               \item If the contribution is a dataset and/or model, the authors should describe the steps taken to make their results reproducible or verifiable.
%               \item Depending on the contribution, reproducibility can be accomplished in various ways. For example, if the contribution is a novel architecture, describing the architecture fully might suffice, or if the contribution is a specific model and empirical evaluation, it may be necessary to either make it possible for others to replicate the model with the same dataset, or provide access to the model. In general. releasing code and data is often one good way to accomplish this, but reproducibility can also be provided via detailed instructions for how to replicate the results, access to a hosted model (e.g., in the case of a large language model), releasing of a model checkpoint, or other means that are appropriate to the research performed.
%               \item While NeurIPS does not require releasing code, the conference does require all submissions to provide some reasonable avenue for reproducibility, which may depend on the nature of the contribution. For example
%                     \begin{enumerate}
%                         \item If the contribution is primarily a new algorithm, the paper should make it clear how to reproduce that algorithm.
%                         \item If the contribution is primarily a new model architecture, the paper should describe the architecture clearly and fully.
%                         \item If the contribution is a new model (e.g., a large language model), then there should either be a way to access this model for reproducing the results or a way to reproduce the model (e.g., with an open-source dataset or instructions for how to construct the dataset).
%                         \item We recognize that reproducibility may be tricky in some cases, in which case authors are welcome to describe the particular way they provide for reproducibility. In the case of closed-source models, it may be that access to the model is limited in some way (e.g., to registered users), but it should be possible for other researchers to have some path to reproducing or verifying the results.
%                     \end{enumerate}
%           \end{itemize}


%     \item {\bf Open access to data and code}
%     \item[] Question: Does the paper provide open access to the data and code, with sufficient instructions to faithfully reproduce the main experimental results, as described in supplemental material?
%     \item[] Answer: \answerTODO{} % Replace by \answerYes{}, \answerNo{}, or \answerNA{}.
%     \item[] Justification: \justificationTODO{}
%     \item[] Guidelines:
%           \begin{itemize}
%               \item The answer NA means that paper does not include experiments requiring code.
%               \item Please see the NeurIPS code and data submission guidelines (\url{https://nips.cc/public/guides/CodeSubmissionPolicy}) for more details.
%               \item While we encourage the release of code and data, we understand that this might not be possible, so “No” is an acceptable answer. Papers cannot be rejected simply for not including code, unless this is central to the contribution (e.g., for a new open-source benchmark).
%               \item The instructions should contain the exact command and environment needed to run to reproduce the results. See the NeurIPS code and data submission guidelines (\url{https://nips.cc/public/guides/CodeSubmissionPolicy}) for more details.
%               \item The authors should provide instructions on data access and preparation, including how to access the raw data, preprocessed data, intermediate data, and generated data, etc.
%               \item The authors should provide scripts to reproduce all experimental results for the new proposed method and baselines. If only a subset of experiments are reproducible, they should state which ones are omitted from the script and why.
%               \item At submission time, to preserve anonymity, the authors should release anonymized versions (if applicable).
%               \item Providing as much information as possible in supplemental material (appended to the paper) is recommended, but including URLs to data and code is permitted.
%           \end{itemize}


%     \item {\bf Experimental setting/details}
%     \item[] Question: Does the paper specify all the training and test details (e.g., data splits, hyperparameters, how they were chosen, type of optimizer, etc.) necessary to understand the results?
%     \item[] Answer: \answerTODO{} % Replace by \answerYes{}, \answerNo{}, or \answerNA{}.
%     \item[] Justification: \justificationTODO{}
%     \item[] Guidelines:
%           \begin{itemize}
%               \item The answer NA means that the paper does not include experiments.
%               \item The experimental setting should be presented in the core of the paper to a level of detail that is necessary to appreciate the results and make sense of them.
%               \item The full details can be provided either with the code, in appendix, or as supplemental material.
%           \end{itemize}

%     \item {\bf Experiment statistical significance}
%     \item[] Question: Does the paper report error bars suitably and correctly defined or other appropriate information about the statistical significance of the experiments?
%     \item[] Answer: \answerTODO{} % Replace by \answerYes{}, \answerNo{}, or \answerNA{}.
%     \item[] Justification: \justificationTODO{}
%     \item[] Guidelines:
%           \begin{itemize}
%               \item The answer NA means that the paper does not include experiments.
%               \item The authors should answer "Yes" if the results are accompanied by error bars, confidence intervals, or statistical significance tests, at least for the experiments that support the main claims of the paper.
%               \item The factors of variability that the error bars are capturing should be clearly stated (for example, train/test split, initialization, random drawing of some parameter, or overall run with given experimental conditions).
%               \item The method for calculating the error bars should be explained (closed form formula, call to a library function, bootstrap, etc.)
%               \item The assumptions made should be given (e.g., Normally distributed errors).
%               \item It should be clear whether the error bar is the standard deviation or the standard error of the mean.
%               \item It is OK to report 1-sigma error bars, but one should state it. The authors should preferably report a 2-sigma error bar than state that they have a 96\% CI, if the hypothesis of Normality of errors is not verified.
%               \item For asymmetric distributions, the authors should be careful not to show in tables or figures symmetric error bars that would yield results that are out of range (e.g. negative error rates).
%               \item If error bars are reported in tables or plots, The authors should explain in the text how they were calculated and reference the corresponding figures or tables in the text.
%           \end{itemize}

%     \item {\bf Experiments compute resources}
%     \item[] Question: For each experiment, does the paper provide sufficient information on the computer resources (type of compute workers, memory, time of execution) needed to reproduce the experiments?
%     \item[] Answer: \answerTODO{} % Replace by \answerYes{}, \answerNo{}, or \answerNA{}.
%     \item[] Justification: \justificationTODO{}
%     \item[] Guidelines:
%           \begin{itemize}
%               \item The answer NA means that the paper does not include experiments.
%               \item The paper should indicate the type of compute workers CPU or GPU, internal cluster, or cloud provider, including relevant memory and storage.
%               \item The paper should provide the amount of compute required for each of the individual experimental runs as well as estimate the total compute.
%               \item The paper should disclose whether the full research project required more compute than the experiments reported in the paper (e.g., preliminary or failed experiments that didn't make it into the paper).
%           \end{itemize}

%     \item {\bf Code of ethics}
%     \item[] Question: Does the research conducted in the paper conform, in every respect, with the NeurIPS Code of Ethics \url{https://neurips.cc/public/EthicsGuidelines}?
%     \item[] Answer: \answerTODO{} % Replace by \answerYes{}, \answerNo{}, or \answerNA{}.
%     \item[] Justification: \justificationTODO{}
%     \item[] Guidelines:
%           \begin{itemize}
%               \item The answer NA means that the authors have not reviewed the NeurIPS Code of Ethics.
%               \item If the authors answer No, they should explain the special circumstances that require a deviation from the Code of Ethics.
%               \item The authors should make sure to preserve anonymity (e.g., if there is a special consideration due to laws or regulations in their jurisdiction).
%           \end{itemize}


%     \item {\bf Broader impacts}
%     \item[] Question: Does the paper discuss both potential positive societal impacts and negative societal impacts of the work performed?
%     \item[] Answer: \answerTODO{} % Replace by \answerYes{}, \answerNo{}, or \answerNA{}.
%     \item[] Justification: \justificationTODO{}
%     \item[] Guidelines:
%           \begin{itemize}
%               \item The answer NA means that there is no societal impact of the work performed.
%               \item If the authors answer NA or No, they should explain why their work has no societal impact or why the paper does not address societal impact.
%               \item Examples of negative societal impacts include potential malicious or unintended uses (e.g., disinformation, generating fake profiles, surveillance), fairness considerations (e.g., deployment of technologies that could make decisions that unfairly impact specific groups), privacy considerations, and security considerations.
%               \item The conference expects that many papers will be foundational research and not tied to particular applications, let alone deployments. However, if there is a direct path to any negative applications, the authors should point it out. For example, it is legitimate to point out that an improvement in the quality of generative models could be used to generate deepfakes for disinformation. On the other hand, it is not needed to point out that a generic algorithm for optimizing neural networks could enable people to train models that generate Deepfakes faster.
%               \item The authors should consider possible harms that could arise when the technology is being used as intended and functioning correctly, harms that could arise when the technology is being used as intended but gives incorrect results, and harms following from (intentional or unintentional) misuse of the technology.
%               \item If there are negative societal impacts, the authors could also discuss possible mitigation strategies (e.g., gated release of models, providing defenses in addition to attacks, mechanisms for monitoring misuse, mechanisms to monitor how a system learns from feedback over time, improving the efficiency and accessibility of ML).
%           \end{itemize}

%     \item {\bf Safeguards}
%     \item[] Question: Does the paper describe safeguards that have been put in place for responsible release of data or models that have a high risk for misuse (e.g., pretrained language models, image generators, or scraped datasets)?
%     \item[] Answer: \answerTODO{} % Replace by \answerYes{}, \answerNo{}, or \answerNA{}.
%     \item[] Justification: \justificationTODO{}
%     \item[] Guidelines:
%           \begin{itemize}
%               \item The answer NA means that the paper poses no such risks.
%               \item Released models that have a high risk for misuse or dual-use should be released with necessary safeguards to allow for controlled use of the model, for example by requiring that users adhere to usage guidelines or restrictions to access the model or implementing safety filters.
%               \item Datasets that have been scraped from the Internet could pose safety risks. The authors should describe how they avoided releasing unsafe images.
%               \item We recognize that providing effective safeguards is challenging, and many papers do not require this, but we encourage authors to take this into account and make a best faith effort.
%           \end{itemize}

%     \item {\bf Licenses for existing assets}
%     \item[] Question: Are the creators or original owners of assets (e.g., code, data, models), used in the paper, properly credited and are the license and terms of use explicitly mentioned and properly respected?
%     \item[] Answer: \answerTODO{} % Replace by \answerYes{}, \answerNo{}, or \answerNA{}.
%     \item[] Justification: \justificationTODO{}
%     \item[] Guidelines:
%           \begin{itemize}
%               \item The answer NA means that the paper does not use existing assets.
%               \item The authors should cite the original paper that produced the code package or dataset.
%               \item The authors should state which version of the asset is used and, if possible, include a URL.
%               \item The name of the license (e.g., CC-BY 4.0) should be included for each asset.
%               \item For scraped data from a particular source (e.g., website), the copyright and terms of service of that source should be provided.
%               \item If assets are released, the license, copyright information, and terms of use in the package should be provided. For popular datasets, \url{paperswithcode.com/datasets} has curated licenses for some datasets. Their licensing guide can help determine the license of a dataset.
%               \item For existing datasets that are re-packaged, both the original license and the license of the derived asset (if it has changed) should be provided.
%               \item If this information is not available online, the authors are encouraged to reach out to the asset's creators.
%           \end{itemize}

%     \item {\bf New assets}
%     \item[] Question: Are new assets introduced in the paper well documented and is the documentation provided alongside the assets?
%     \item[] Answer: \answerTODO{} % Replace by \answerYes{}, \answerNo{}, or \answerNA{}.
%     \item[] Justification: \justificationTODO{}
%     \item[] Guidelines:
%           \begin{itemize}
%               \item The answer NA means that the paper does not release new assets.
%               \item Researchers should communicate the details of the dataset/code/model as part of their submissions via structured templates. This includes details about training, license, limitations, etc.
%               \item The paper should discuss whether and how consent was obtained from people whose asset is used.
%               \item At submission time, remember to anonymize your assets (if applicable). You can either create an anonymized URL or include an anonymized zip file.
%           \end{itemize}

%     \item {\bf Crowdsourcing and research with human subjects}
%     \item[] Question: For crowdsourcing experiments and research with human subjects, does the paper include the full text of instructions given to participants and screenshots, if applicable, as well as details about compensation (if any)?
%     \item[] Answer: \answerTODO{} % Replace by \answerYes{}, \answerNo{}, or \answerNA{}.
%     \item[] Justification: \justificationTODO{}
%     \item[] Guidelines:
%           \begin{itemize}
%               \item The answer NA means that the paper does not involve crowdsourcing nor research with human subjects.
%               \item Including this information in the supplemental material is fine, but if the main contribution of the paper involves human subjects, then as much detail as possible should be included in the main paper.
%               \item According to the NeurIPS Code of Ethics, workers involved in data collection, curation, or other labor should be paid at least the minimum wage in the country of the data collector.
%           \end{itemize}

%     \item {\bf Institutional review board (IRB) approvals or equivalent for research with human subjects}
%     \item[] Question: Does the paper describe potential risks incurred by study participants, whether such risks were disclosed to the subjects, and whether Institutional Review Board (IRB) approvals (or an equivalent approval/review based on the requirements of your country or institution) were obtained?
%     \item[] Answer: \answerTODO{} % Replace by \answerYes{}, \answerNo{}, or \answerNA{}.
%     \item[] Justification: \justificationTODO{}
%     \item[] Guidelines:
%           \begin{itemize}
%               \item The answer NA means that the paper does not involve crowdsourcing nor research with human subjects.
%               \item Depending on the country in which research is conducted, IRB approval (or equivalent) may be required for any human subjects research. If you obtained IRB approval, you should clearly state this in the paper.
%               \item We recognize that the procedures for this may vary significantly between institutions and locations, and we expect authors to adhere to the NeurIPS Code of Ethics and the guidelines for their institution.
%               \item For initial submissions, do not include any information that would break anonymity (if applicable), such as the institution conducting the review.
%           \end{itemize}

%     \item {\bf Declaration of LLM usage}
%     \item[] Question: Does the paper describe the usage of LLMs if it is an important, original, or non-standard component of the core methods in this research? Note that if the LLM is used only for writing, editing, or formatting purposes and does not impact the core methodology, scientific rigorousness, or originality of the research, declaration is not required.
%           %this research? 
%     \item[] Answer: \answerTODO{} % Replace by \answerYes{}, \answerNo{}, or \answerNA{}.
%     \item[] Justification: \justificationTODO{}
%     \item[] Guidelines:
%           \begin{itemize}
%               \item The answer NA means that the core method development in this research does not involve LLMs as any important, original, or non-standard components.
%               \item Please refer to our LLM policy (\url{https://neurips.cc/Conferences/2025/LLM}) for what should or should not be described.
%           \end{itemize}

% \end{enumerate}


\end{document}